\usemodule[memos]
\usemodule[basicexam][mode={teacher,check}]

\setupinteraction[state=start]
\setuptextrules[style=\tf\ss,before=,after=,rulethickness=.5pt]
\setuphead[section][style=\ss\bfb,color=black]
\setuphead[subsection][style=\ss\bf,color=blue]
\setuptyping[style=\ttxx\setuplocalinterlinespace[line=2ex],
             option={context},
             before={\vbox{\hrule height .5pt\par
                     \textrule{EXAMPLE}%
                     \hrule height .5pt\vrule width0pt depth 4pt}\par},
              after={\vbox{\hrule height .5pt\strut}},
             ]

\starttext
\dontcomplain

\placecontent
\chapter{Overview}

% ==============================================
% Module Overview
% ==============================================

\section{Introduction to basicexam Module}

The basicexam module is a ConTeXt module designed for creating educational
materials such as exams, quizzes, and practice tests. It provides a flexible
framework for defining questions, choices, answers, and scoring.

\subsection{Module Features}

\startitemize[packed]
\item Multiple question types including choice questions, fill-in-the-blank, and open-ended
\item Support for teacher/student modes with answer visibility control
\item Two-up layout option for printing exams on both sides of paper
\item Customizable styling for questions, answers, and scoring
\item Inline choice questions for compact presentation
\item Script areas for student answers
\item Comprehensive point tracking and scoring
\stopitemize

\subsection{Module Parameters}

The module accepts the following parameters when loaded:

\starttabulate[|l|l|p|]
\HL
\NC Parameter \NC Values \NC Description \NC\NR
\HL
\NC mode      \NC        \NC  Controls the display mode:\NC\NR
\NC  \NC - student \NC Shows questions only\NC\NR
\NC  \NC - teacher \NC Shows questions with answers\NC\NR
\NC  \NC - check   \NC Debug mode with console output \NC\NR
\HL
\NC layout \NC \NC Layout configuration:\NC\NR
\NC \NC - twoup \NC Two virtual pages per physical page\NC\NR
\NC \NC - none  \NC Default single-page layout \NC\NR
\HL
\NC footer \NC \NC Footer configuration:\NC\NR
\NC \NC - default \NC Shows subject and page number\NC\NR
\NC \NC - none    \NC No footer \NC\NR
\HL
\stoptabulate

\setupquestion[answer={default answer},point=1]
\setupproblem [left=(,right=),stopper=,width=\widthofstring{(00)}]

\section{Choice Questions}

This section demonstrates various implementations of choice questions,
showing both verbose and concise syntax options.

\subsection{Basic Choice Question Syntax}

The following examples show the two primary syntax formats for creating choice questions:

\startitemize[packed]
\item Verbose format using \type{\startquestion}/\type{\stopquestion} environments
\item Quick format using the \type{\question} macro with comma-separated choices
\stopitemize

Additionally, they demonstrate column control through the optional parameter
(8 columns in the second example) and integration with external figures.

\startbuffer
\startquestion
  This is the question stem of the selection question.
  \startchoice
    \startcitem Choice One   \stopcitem
    \startcitem Choice Two   \stopcitem
    \startcitem Choice Three \stopcitem
    \startcitem Choice Four  \stopcitem
  \stopchoice
\stopquestion

\question{This is the question stem of the selection question.
   \choice{Choice One,Choice Two,Choice Three,Choice Four}
}

\start
\startquestion
  This is the question stem of the selection question.
  \startchoice[8]
    \startcitem Choice One   \stopcitem
    \startcitem Choice Two   \stopcitem
    \startcitem Choice Three \stopcitem
    \startcitem Choice Four  \stopcitem
  \stopchoice
\stopquestion

\question{This is the question stem of the selection question.
   \choice[8]{Choice One,Choice Two,Choice Three,Choice Four}
}

\startplacefigure [location={none,right,2*hang}]
  \externalfigure [][height=5\baselineskip]
\stopplacefigure
\question
{\input knuthmath
\choice[one,vertical]{Choice One,Choice Two,Choice Three,Choice Four}
}
\stop
\stopbuffer

\typebuffer

\getbuffer

\subsection{Questions with Answers and Points}

This section demonstrates how to define questions with:

\startitemize[packed]
\item Explicit answer marking using \type{[*]} syntax
\item Point values and point display
\item Detailed answer explanations
\item Reference labels for cross-referencing
\stopitemize

The answer environment supports multiple formatting options and will automatically
show/hide based on the module's mode parameter (teacher/student).

\startbuffer
\startquestion[reference=Q,
               showanswer=true,
               showpoint=true,
               point=4,
               answer=Q]
  This is the question stem of the selection question \fillin{Answer Q}.
  \startchoice
    \startcitem    Choice One   \stopcitem
    \startcitem    Choice Two   \stopcitem
    \startcitem[*] Choice Three \stopcitem
    \startcitem    Choice Four  \stopcitem
  \stopchoice
  \startanswer
    This is a detailed explanation of this topic.
    It should be noted that the answer will be overwritten by the next answer. 
    Three answers are set for this question, but only the last answer will be selected in the end.
  \stopanswer
\stopquestion

\question[reference=QQ,
          showanswer=true,
          showpoint=true,
          point=4,
          answer=QQ]
  {This is the question stem of the selection question.\fillin{Answer QQ}.
  \choice{Choice One,Choice Two,Choice Three,{[*]Choice Four}}
  \startanswer
  This is a detailed explanation of this topic.
  It should be noted that the answer will be overwritten by the next answer. 
  Three answers are set for this question, but only the last answer will be selected in the end.
  \stopanswer
  }

\stopbuffer

\typebuffer

\getbuffer

\subsection{Answer Script Areas}

This section shows how to add script areas where students can write responses.
Key features demonstrated:

\startitemize[packed]
\item \type{showscript=true} to enable answer spaces
\item Custom script content with light gray prompt text
\item Frame control for answer areas
\item Vertical spacing adjustments for handwritten answers
\stopitemize

Script areas are particularly useful for open-ended questions requiring longer responses.

\startbuffer
\startquestion[reference=Q,
               showpoint=true,
               point=4,
               answer=Q]
  This is the question stem of the selection question \fillin{Answer Q}.
  \startchoice
    \startcitem    Choice One   \stopcitem
    \startcitem    Choice Two   \stopcitem
    \startcitem[*] Choice Three \stopcitem
    \startcitem    Choice Four  \stopcitem
  \stopchoice
  \startanswer[showscript=true,showanswer=false,
               scriptcontent={\centerline{\lightgray\ss WRITE YOUR ANSWER HERE.}\par
                              \null\vskip2\baselineskip\relax}]
    This is a detailed explanation of this topic.
    It should be noted that the answer will be overwritten by the next answer. 
    Three answers are set for this question, but only the last answer will be selected in the end.
  \stopanswer
\stopquestion

\question[showpoint=true,point=4,]
  {This is the question stem of the selection question.\relax
   \choice{Choice One,Choice Two,Choice Three,Choice Four}}
\startanswer[showscript=true,frame=off,scriptcontent=\null\blank[3*line]]

\stopanswer

\question[showpoint=true,point=4,]
  {This is the question stem of the selection question.\relax
   \choice{Choice One,Choice Two,Choice Three,Choice Four}}
\stopbuffer

\typebuffer

\getbuffer

\subsection{Styling Customization}

This section demonstrates the extensive styling options available:

\startitemize[packed]
\item Question text styling (font family, size, color)
\item Choice option formatting
\item Point display customization
\item Answer formatting and color coding
\item Label styling and positioning
\stopitemize

These styling parameters can be applied globally through \type{\setupquestion}
or locally for individual questions.

\startbuffer
\startquestion[style=ss,pointlabel={total , pts},point=2,showpoint=true]
  This is the question stem of the selection question.
  \startchoice[style=tt,color=red]
    \startcitem Choice One   \stopcitem
    \startcitem Choice Two   \stopcitem
    \startcitem Choice Three \stopcitem
    \startcitem Choice Four  \stopcitem
  \stopchoice
\stopquestion

\startquestion[answerstyle=tt,answercolor=darkred,
               point=2,
               pointstyle=tt, pointcolor=darkred,
               textstyle=ss,textcolor=blue,
               showpoint=true,showanswer=true]
  This is the question stem of the selection question.
  \startchoice[answerstyle=\tt,answercolor=green,textstyle=rm,textcolor=darkred,]
    \startcitem Choice One      \stopcitem
    \startcitem Choice Two      \stopcitem
    \startcitem[*] Choice Three \stopcitem
    \startcitem Choice Four     \stopcitem
  \stopchoice
\stopquestion
\stopbuffer

\typebuffer

\getbuffer

\subsection{Inline Choice Questions}

Inline choice questions allow for compact presentation where choices appear
within the flow of the question text rather than as separate list items.
This format is particularly useful for short questions or when space is limited.

The example demonstrates both environment-based (\type{\startinlinechoice}) and macro-based (\type{\inlinechoice}) inline choice syntax.

\startbuffer
\startquestion
  This is the question stem of the selection question.\relax
  \startinlinechoice
    \startcitem Choice One   \stopcitem
    \startcitem Choice Two   \stopcitem
    \startcitem Choice Three \stopcitem
    \startcitem Choice Four  \stopcitem
  \stopinlinechoice
\stopquestion

\question{This is the question stem of the selection question.\relax
   \inlinechoice{Choice One,Choice Two,Choice Three,Choice Four}
}

\startquestion[showanswer=true]
  This is the question stem of the selection question.\relax\break
  \startinlinechoice
    \startcitem[*] Choice One\stopcitem
    \startcitem Choice Two   \stopcitem
    \startcitem Choice Three \stopcitem
    \startcitem Choice Four  \stopcitem
  \stopinlinechoice
\stopquestion

\question[showanswer=true]{This is the question stem of the selection question.\relax\break\hfill
   \inlinechoice{Choice One,{[*]Choice Two},Choice Three,Choice Four}
}
\stopbuffer

\typebuffer

\getbuffer

\startbuffer
\startquestion
  This is the question stem of the selection question.\relax
  \startchoice[vertical,text][]
    \startcitem Choice One   \stopcitem
    \startcitem[*] Choice Two   \stopcitem
    \startcitem Choice Three \stopcitem
    \startcitem Choice Four  \stopcitem
  \stopchoice
\stopquestion

\question[showanswer=true]{This is the question stem of the selection question.\relax
   \choice[vertical,text][]{Choice One,{[*]Choice Two},Choice Three,Choice Four}
}

\definechoice[newinlinechoice][inlinechoice=true,correctsymbol=default]
\startquestion
  This is the question stem of the selection question.\relax
  \startnewinlinechoice
    \startcitem Choice One   \stopcitem
    \startcitem[*] Choice Two   \stopcitem
    \startcitem Choice Three \stopcitem
    \startcitem Choice Four  \stopcitem
  \stopnewinlinechoice
\stopquestion

\question[showanswer=true]{This is the question stem of the selection question.\relax
   \newinlinechoice{Choice One,{[*]Choice Two},Choice Three,Choice Four}
}
\stopbuffer

\typebuffer

\getbuffer

\subsection{multichoice}

\startbuffer
\startquestion[showanswer=true]
  This is the question stem of the selection question.\relax
  \startchoice
    \startcitem[*]   Choice One\stopcitem
    \startcitem      Choice Two\stopcitem
    \startcitem[*] Choice Three\stopcitem
    \startcitem     Choice Four\stopcitem
  \stopchoice
\stopquestion

\question[showanswer=true]{This is the question stem of the selection question.\relax
   \choice{{[*]Choice One},{[*]Choice Two},Choice Three,Choice Four}
}

\startquestion[showanswer=true]
  This is the question stem of the selection question.\relax
  \startinlinechoice
    \startcitem[*]   Choice One\stopcitem
    \startcitem      Choice Two\stopcitem
    \startcitem[*] Choice Three\stopcitem
    \startcitem     Choice Four\stopcitem
  \stopinlinechoice
\stopquestion

\question[showanswer=true]{This is the question stem of the selection question.\relax
   \inlinechoice{Choice One,{[*]Choice Two},Choice Three,{[*]Choice Four}}
}
\stopbuffer

\typebuffer

\getbuffer

\subsection{more instance}

\startbuffer
  \definechoice[choiceexample]
  \setupchoice
    [choiceexample]
    [option={text,6,packed,joinedup,nowhite,},
     inlinechoice=true,
     correctsymbol={\symbol{markcirclecheck}},
     inlineleft={!!\nobreak},
     inlineright={\nobreak !!},
     answerstyle=\ss,
     separator={\removeunwantedspaces,\space},]
     
\startquestion[showanswer=true]
  This is the question stem of the selection question.\relax
  \startchoiceexample
    \startcitem[*]   Choice One\stopcitem
    \startcitem      Choice Two\stopcitem
    \startcitem[*] Choice Three\stopcitem
    \startcitem     Choice Four\stopcitem
  \stopchoiceexample
\stopquestion
\stopbuffer

\typebuffer

\getbuffer

\subsection{point calculation}

\startbuffer
\question[point={1.5},showpoint=true]{AXXXA\choice{dummy text,XXX,YYY,ZZZ}}
\question[point={2.2},showpoint=true]{AXXXA\choice{dummy text dummy text,XXX,YYY,ZZZ}}
\question[point={1.8},showpoint=true]{AXXXA\choice{dummy text dummy text dummy text dummy text,XXX,YYY,ZZZ}}
\stopbuffer

\typebuffer

\getbuffer

\section{Question}

\subsection{question \& problem}

\startbuffer
\startquestion[start=1,option={n,packed,joinedup,nowhite},]
This is the main part of the question.
\stopquestion

\question
{This is the main part of the question.}

\startquestion
This is the main part of the question.
  \startquestion
    This is the main part of the question.
      \startquestion
        This is the main part of the question.
      \stopquestion
  \stopquestion
  \startquestion
    This is the main part of the question.
      \startquestion
        This is the main part of the question.
      \stopquestion
      \startquestion
        This is the main part of the question.
      \stopquestion
      \startquestion
        This is the main part of the question.
      \stopquestion
  \stopquestion
\stopquestion

\startquestion
This is the main part of the question.
  \startproblem
  \startpitem  This is the first small question.\stoppitem
  \startpitem This is the second small question.\stoppitem
  \stopproblem
\stopquestion

\question
{This is the main part of the question.
  \problem[textstyle=ss]
    {\pitem{This is the first small question.}
     \pitem{This is the second small question.}
  }
}
\stopbuffer

\typebuffer

\getbuffer

\subsection{True/False questions}

\startbuffer
\startquestion
This is the main part of the question.
  \startproblem
  \startpitem  \fillin{} This is the first  small question.\stoppitem
  \startpitem  \fillin{} This is the second small question.\stoppitem
  \stopproblem
\stopquestion

\start
\setupfillin[empty=none,style=\tta]
\startquestion
This is the main part of the question.
  \startproblem
  \startpitem  \fillin{T} This is the first  small question.\stoppitem
  \startpitem  \fillin{F} This is the second small question.\stoppitem
  \startpitem  This is the first  small question.\hfill\fillin{\symbol{markcheck}}\stoppitem
  \startpitem  This is the second small question.\hfill\fillin{\symbol{markcross}}\stoppitem
  \stopproblem
\stopquestion
\stop
\stopbuffer

\typebuffer

\getbuffer

\subsection{example with answer}

\startbuffer
\setupquestion[showanswer=true,showpoint=true,point=1]
\setupanswer  [showanswer=true,showpoint=true]
\startquestion[start=1,option={n,packed,joinedup,nowhite},answer=aaa]
This is the main part of the question.
\startanswer This is a detailed explanation of this topic.\stopanswer
\stopquestion

\question[answer=bbb]
{This is the main part of the question.
\startanswer This is a detailed explanation of this topic.\stopanswer}

\startquestion
This is the main part of the question.
  \startproblem
  \startpitem[answer=eee]
    This is the first small question.
    \startanswer This is a detailed explanation of this topic.\stopanswer
  \stoppitem
  \startpitem[answer=fff]
    This is the second small question.
    \startanswer This is a detailed explanation of this topic.\stopanswer
  \stoppitem
  \stopproblem
\stopquestion

\question
{This is the main part of the question.
  \problem[textstyle=ss]
    {\pitem[answer=ggg]{This is the first small question.
        \answer{This is a detailed explanation of this topic.}}
     \pitem[answer=hhh]{This is the second small question.
        \answer{This is a detailed explanation of this topic.}}
     \pitem[answer=iii]{This is the second small question.
        \answer{This is a detailed explanation of this topic.}}
  }
}
\stopbuffer

\typebuffer

\start\getbuffer\stop

\setupanswer
  [scriptcontent={
    \centerline{\lightgray
    \ss WRITE YOUR ANSWER HERE.}\par
    \null\vskip2\baselineskip\relax}]

\startbuffer
\startpitem[answer={ccc,will be replace}]
   Lorem ipsum dolor sit amet,
   \fillin{Lorem ipsum,the real answer}
\startanswer
 Here is answer.
\stopanswer
\stoppitem
\stopbuffer

\subsection{fillin }

\start

\setupquestion[answer=cxxx,showanswer=false,point=2,showpoint=false]

\startquestion[showanswer=false]
  Lorem ipsum dolor sit amet,
  \startproblem
    {\setupfillin[empty=number, n=5]\getbuffer}
    {\setupfillin[empty=number, align=autoright,n=5]\getbuffer}
    
    {\setupfillin[empty=yes,    n=5]\getbuffer}
    {\setupfillin[empty=yes,    align=autoright,n=5]\getbuffer}

    {\setupfillin[empty=none,   n=5]\getbuffer}
    {\setupfillin[empty=none,   align=autoright,n=5]\getbuffer}
  \stopproblem
\stopquestion

\setupquestion[answer=cxxx,showanswer=true,point=2,showpoint=true]

\startquestion[showanswer=true]
  Lorem ipsum dolor sit amet,
  \startproblem
    \startpitem[answer={ccc,will be replace},point=5]
               Lorem ipsum dolor sit amet,
               \fillin{Lorem ipsum,the real answer 1}
               \fillin{Lorem ipsum,the real answer 2}\stoppitem
    {\setupfillin[empty=number, n=5]\getbuffer}
    {\setupfillin[empty=yes,    n=5]\getbuffer}
    {\setupfillin[empty=none,   n=5]\getbuffer}
    {\setupfillin[empty=number, align=autoright,n=5]\getbuffer}
    {\setupfillin[empty=yes,    align=autoright,n=5]\getbuffer}
    {\setupfillin[empty=none,   align=autoright,n=5]\getbuffer}
  \stopproblem
\stopquestion

\stop

\section{Material}

\startbuffer
\setupmaterial[title] [style=\ssa]
\setupmaterial[author][style=\tta]
\setupmaterial[source][style=ss]
\setupmaterial[number][numberconversion=I]
\setupmaterial[indicator][numberconversion=A,reset=true]

\startmaterial[title={Knuth 1},author={Mos},source={Yelu}]
  \input knuth\relax
  \indicator{Some Texts 1}
  \indicator{Some Texts 2} 
\stopmaterial

\startquestion[start=1,option={n,packed,joinedup,nowhite}]
  \indicator{Some Texts 1}
  \choice{Some Texts,Some Texts,Some Texts,Some Texts}
\stopquestion
\startquestion
  \indicator{Some Texts 2}
  \choice{Some Texts,Some Texts,Some Texts,Some Texts}
\stopquestion

\startmaterial[title={Knuth 2},author={Mos},source={Yelu}]
\input knuth\relax\indicator{Some Texts 3} \indicator{Some Texts 4}
\stopmaterial
\startquestion
  \indicator{Some Texts 3}
  \choice{Some Texts,Some Texts,Some Texts,Some Texts}
\stopquestion
\startquestion
  \indicator{Some Texts 4}
  \choice{Some Texts,Some Texts,Some Texts,Some Texts}
\stopquestion
\stopbuffer

\typebuffer
\getbuffer

\section{Close}

\startclose[showanswer=true,showpoint=true,point=5,answer=nil]
On Oct. 11, hundreds of runners competed in a cross-country race in Minnesota.
Melanie Bailey should have \closechoice[designed,followed,changed,finished] the
course earlier than she did. Her \closechoice[delay,chance,trouble,excuse] came
because she was carrying a \closechoice[judge,volunteer,classmate,competitor]
across the finish line.

As reported by a local newspaper, Bailey was more than two-thirds of the way
through her \closechoice[race,school,town,training] when a runner in front of
her began crying in pain. She \closechoice[agreed,returned,stopped,promised]
to help her fellow runner, Danielle Lenoue. Bailey took her arm to see if she
could walk forward with \closechoice[courage,aid,patience,advice] . She couldn't.
Bailey then \closechoice[went away,stood up,stepped aside,bent down] to let Lenoue
climb onto her back and carried her all the way to the finish line, then another
300 feet to where Lenoue could get \closechoice[medical,public,constant,equal] attention.

Once there, Lenoue was \closechoice[interrupted,assessed,identified,appreciated]
and later taken to a hospital, where she learned that she had serious injuries in
one of her knees. She would have struggled with extreme \closechoice[hunger,pain,cold,tiredness]
to make it to that aid checkpoint without Bailey's help.

As for Bailey, she is more \closechoice[worried,ashamed,confused,discouraged]
about why her act is considered a big \closechoice[game,problem,lesson,deal] .
"She was just crying. I couldn't \closechoice[leave,cure,bother,understand] her,"
Bailey told the reporter. "I feel like I was just doing the right thing.”

Although the two young women were strangers before the \closechoice[ride,test,meet,show] ,
they've since become friends. Neither won the race, but the \closechoice[secret,display,benefit,exchange]
of human kindness won the day.
\stopclose


\section{Chinese Poem}

\setupverse [sample={滾滾長江東逝水浪花淘盡英雄},skip=left,before=\blank,after=\blank]
\setupverse [title] [align=center,style=bf,skip=both]
\setupverse [author][align=flushright,style=it,skip=both]
\setupverse [source][align=flushright,style=it,skip=both]

\startbuffer
\startverse[title=临江仙 $\cdot$ 滚滚长江东逝水,author=明 $\cdot$ 杨慎]
滾滾長江東逝水浪花淘盡英雄\\
是非成敗轉頭空\\
青山依舊在 幾度夕陽紅\\
白髮漁樵江渚上 慣看秋月春風\\
一壺濁酒喜相逢\\
古今多少事 都付笑談中
\stopverse
\stopbuffer

\typebuffer

\getbuffer

\startbuffer
\setupverse[skip=none,align=center]
\stopbuffer

\typebuffer

\getbuffer

\section{Lexicalentry}

\startbuffer
\startlexicalentry[class=名,accent=1]{太陽}
  \startlexicallists
    \item 太陽系の中心をなし、地球にもっとも近い恒星。お日さま。\antonym{太陰}。
    \item 光明・希望の象徴。「心の\alias」
  \stoplexicallists
\stoplexicalentry
\stopbuffer
\typebuffer
\getbuffer

\section{Answer Output (typeanswer)}

The \type{\typeanswer} command is used to output collected answers, supporting multiple output styles and customizable parameters.

\subsection{Core Parameters}

\starttabulate[|l|p|]
\HL
\NC Parameter \NC Description \NC\NR
\HL
\NC \type{start} \NC Starting question number \NC\NR
\NC \type{stop} \NC Ending question number \NC\NR
\NC \type{total} \NC Total number of questions (alternative to stop) \NC\NR
\NC \type{chapter} \NC Chapter number (default: current chapter) \NC\NR
\NC \type{alternative} \NC Output style \NC\NR
\NC \type{rows}, \type{columns} \NC Table rows and columns \NC\NR
\NC \type{separator} \NC Separator (for text style) \NC\NR
\NC \type{left}, \type{right} \NC Content to wrap around answers \NC\NR
\NC \type{before}, \type{after} \NC Content before and after each item \NC\NR
\HL
\stoptabulate

\subsection{Predefined Alternatives}

\startbuffer
% List format (default)
\typeanswer[start=1,total=5,alternative=list]{}

% Table format
\typeanswer[start=1,total=10,alternative=table,rows=2,columns=5]{}

% Text format
\typeanswer[start=1,total=5,alternative=text,separator={\quad},left={\(},right={\)}]{}

% Grid format
\typeanswer[start=1,total=10,alternative=grid,rows=5,columns=2]{}

% Direct access
\typeanswer[alternative=direct]{1-Q-5}
\stopbuffer

\typebuffer

\getbuffer

\subsection{Custom Alternative}

Users can define custom output styles using \type{\definetypeansweralternative}:

\startbuffer
\unprotect
\definetypeansweralternative
  [mycustom]
  [renderingsetup=\????typeanswerrenderings:mycustom]

\startsetups[\????typeanswerrenderings:mycustom]
  \scratchcounter=\numexpr\c_typeanswer_stop-\c_typeanswer_start+1\relax
  \c_typeanswer_index=\c_typeanswer_start
  \dorecurse\scratchcounter{%
    \bf Question \typeanswergetnumber: \typeanswergetanswer (\typeanswergetpoint points)\par
    \advance\c_typeanswer_index by 1\relax}
\stopsetups
\protect

\typeanswer[start=1,total=3,alternative=mycustom]{}
\stopbuffer

\typebuffer

\getbuffer

\subsection{User Commands}

User commands available for customizing alternatives:

\starttabulate[|l|p|]
\HL
\NC Command \NC Description \NC\NR
\HL
\NC \type{\typeanswercurrentindex} \NC Current index \NC\NR
\NC \type{\typeanswergetnumber} \NC Get question number \NC\NR
\NC \type{\typeanswergetanswer} \NC Get answer \NC\NR
\NC \type{\typeanswergetpoint} \NC Get point value \NC\NR
\NC \type{\typeanswergetdirect{id}{field}} \NC Directly get specified data \NC\NR
\HL
\stoptabulate

\section{Score}

\startbuffer
\startanswer[showanswer=true]
 some text \score{2}
 some text \score{4}
 \resetscore
 \setupscore[score=calculation]
 some text \score[type=cdotfill]{2}
 some text \score[type=space]{4}
 some text \score[type=hfill]{3}
 \resetscore
 \setupscore[score=complex]
 some text \score{2}
 some text \score{4}
 some text \score{3}
\stopanswer
\stopbuffer

\typebuffer

\getbuffer

\section{Paper Title}

The \type{\definepapertitle}, \type{\setuppapertitle}, and \type{\makepapertitle}
commands are used to create professional test paper headers with customizable
information sections.

\startitemize[packed]
\item \type{\definepapertitle[PaperTitleName]} --- Define a new paper title instance
\item \type{\setuppapertitle[PaperTitleName][...]} --- Configure paper title settings
\item \type{\makepapertitle[PaperTitleName]} --- Render the paper title
\stopitemize

\subsection{Paper Title Parameters}

The paper title system supports the following key parameters:

\starttabulate[|l|p|]
\HL
\NC Parameter \NC Description \NC\NR
\HL
\NC list \NC Order of elements to display (default: secret, title, subject, information, signinformation, notice) \NC\NR
\NC secret \NC Secret/classification text (e.g., "绝密 ★ 启用前") \NC\NR
\NC title \NC Main paper title \NC\NR
\NC subject \NC Subject name (appended with subjectsuffix) \NC\NR
\NC subjectsuffix \NC Suffix for subject (default: 学科试卷) \NC\NR
\NC information \NC General information (time, total score, etc.) \NC\NR
\NC signinformation \NC Student sign-in information fields \NC\NR
\NC notice \NC Important notices/instructions \NC\NR
\HL
\stoptabulate

Each element in the list automatically generates style-related parameters:
\startitemize[packed]
\item \type{<element>style} --- Text style
\item \type{<element>color} --- Text color
\item \type{<element>align} --- Text alignment
\item \type{before<element>} --- Content before the element
\item \type{after<element>} --- Content after the element
\stopitemize

\startbuffer
\definepapertitle[myexam]
\setuppapertitle[myexam][
    list={secret,examtitle,subject,information,notice},
    subjectsuffix=学科试卷,
    secretstyle=\ss,
    examtitlestyle=\ssa,
    subjectstyle=\ssb,
    informationstyle=\ttx,
    noticestyle=\rm\it,
    secretalign=flushleft,
    examtitlealign=center,
    subjectalign=center,
    informationalign=center,
    noticealign=flushleft,
    secret={绝密 ★ 启用前},
    examtitle={2024 年普通高等学校招生全国统一考试},
    subject={英语},
    information={总分:150 分，考试时间:120 分钟},
    notice={注意事项：
      \startitemize[n,packed,joinedup]
        \item 答题前，务必将自己的姓名、准考证号填写在答题卡规定的位置上。
        \item 答选择题时，必须使用 2B 铅笔将答题卡上对应题目的答案标号涂黑。
        \item 考试结束后，将试题卷和答题卡一并交回。
      \stopitemize},
]
\makepapertitle[myexam]
\stopbuffer

\typebuffer

\start\getbuffer\stop

\section{Writing Environment}

The \type{writing} environment is designed for composition/writing questions.
It supports both standalone writing tasks and nested sub-writing tasks.

\starttabulate[|l|l|p|]
\HL
\NC Parameter \NC Values \NC Description \NC\NR
\HL
\NC point \NC number \NC Point value for the writing task \NC\NR
\NC answer \NC text \NC Reference answer text \NC\NR
\NC showanswer \NC true/false \NC Show/hide reference answer \NC\NR
\NC showpoint \NC true/false \NC Show/hide point value \NC\NR
\NC totalanswer \NC text \NC Override total answer display \NC\NR
\HL
\stoptabulate

\subsection{Basic Writing Task}

\startbuffer
\setupwriting[showanswer=true,showpoint=true]
\startwriting[point=20]
  Write an essay about your favorite season. Your essay should be at least 200 words.
  \startanswer
    This is a reference composition that demonstrates the expected quality
    and length of the student's response. It covers the main points and
    shows good organization and language use.
  \stopanswer
\stopwriting
\stopbuffer

\typebuffer

\getbuffer

\subsection{Nested Writing Tasks}

When a writing task contains multiple sub-topics, use \type{subwriting}:

\startbuffer
\setupwriting[showanswer=true,showpoint=true]
\startwriting
  Choose one of the following topics and write an essay:
  \startsubwriting[point=15]
    Topic A: Describe a memorable travel experience.
    \startanswer
      Reference essay for Topic A...
    \stopanswer
  \stopsubwriting
  \startsubwriting[point=15]
    Topic B: Discuss the importance of environmental protection.
    \startanswer
      Reference essay for Topic B...
    \stopanswer
  \stopsubwriting
\stopwriting
\stopbuffer

\typebuffer

\getbuffer

\section{Optclose Environment}

The \type{optclose} environment is a variant of the \type{close} environment,
designed for fill-in-the-blank questions where the answer is provided inline.
Use \type{\ocitem} to mark each blank.

\starttabulate[|l|l|p|]
\HL
\NC Parameter \NC Values \NC Description \NC\NR
\HL
\NC answer \NC text \NC The correct answer for this blank \NC\NR
\NC point \NC number \NC Point value (default: 1) \NC\NR
\NC reset \NC true/false \NC Reset numbering for each environment \NC\NR
\NC continue \NC true/false \NC Continue from previous question number \NC\NR
\HL
\stoptabulate

\startbuffer
\startoptclose[showanswer=true,showpoint=true,point=4]
The latest \ocitem[answer=engineering,point=6]{engineer} techniques are applied to create
this protective \ocitem[answer=functional]{function} structure that is also
beautiful. The design features ten steel "sepals" made of glass and aluminium.
These sepals open on warm days \ocitem[answer=to give]{give} the inside plants
sunshine and fresh air. In cold weather, the structure stays \ocitem[answer=closed]{close}
to protect the plants.
\stopoptclose
\stopbuffer

\typebuffer

\getbuffer

\section{Writing Box and Grid}

The \type{\writingbox} and \type{\writinggrid} commands create formatted areas
for handwritten answers, commonly used in composition sections.

\subsection{Writing Box}

\starttabulate[|l|l|p|]
\HL
\NC Parameter \NC Values \NC Description \NC\NR
\HL
\NC column \NC number \NC Number of columns \NC\NR
\NC row \NC number \NC Number of rows \NC\NR
\NC evenheight \NC dimension \NC Height of even rows \NC\NR
\NC oddheight \NC dimension \NC Height of odd rows \NC\NR
\NC evencolor \NC color \NC Color of even rows \NC\NR
\NC oddcolor \NC color \NC Color of odd rows \NC\NR
\HL
\stoptabulate

\startbuffer
\writingbox
  [column=5,
   row=5,
   evencolor=blue,
   oddcolor=gray,
   evenheight=1em,
   oddheight=.5em,
   height=5cm,
   width=\textwidth,]{}
\stopbuffer

\typebuffer

\getbuffer

\subsection{Writing Grid}

\starttabulate[|l|l|p|]
\HL
\NC Parameter \NC Values \NC Description \NC\NR
\HL
\NC column \NC number \NC Number of columns \NC\NR
\NC row \NC number \NC Number of rows \NC\NR
\NC cellwidth \NC dimension \NC Width of each cell \NC\NR
\NC cellcolor \NC color \NC Color of grid lines \NC\NR
\HL
\stoptabulate

\startbuffer
\writinggrid
  [cellwidth=1cm,
   column=10,
   row=5,
   cellcolor=darkgray,
   height=5cm,
   width=\textwidth,]{}
\stopbuffer

\typebuffer

\getbuffer

\section{Symbols}

The module provides a comprehensive set of marking symbols for answers.
These symbols require the {\bf NotoSansSymbols2} font to be installed.

\starttabulate[|l|r|l|r|]
\HL
\NC Command \NC Symbol \NC Command \NC Symbol \NC\NR
\HL
\NC \type{\symbol{marksquare}} \NC \symbol{marksquare} \NC \type{\symbol{markcircle}} \NC \symbol{markcircle} \NC\NR
\NC \type{\symbol{markcheck}} \NC \symbol{markcheck} \NC \type{\symbol{markcross}} \NC \symbol{markcross} \NC\NR
\NC \type{\symbol{markheavycheck}} \NC \symbol{markheavycheck} \NC \type{\symbol{markheavycross}} \NC \symbol{markheavycross} \NC\NR
\NC \type{\symbol{marksquarecheck}} \NC \symbol{marksquarecheck} \NC \type{\symbol{marksquarecross}} \NC \symbol{marksquarecross} \NC\NR
\NC \type{\symbol{marksquareheavycheck}} \NC \symbol{marksquareheavycheck} \NC \type{\symbol{marksquareheavycross}} \NC \symbol{marksquareheavycross} \NC\NR
\NC \type{\symbol{markcirclecheck}} \NC \symbol{markcirclecheck} \NC \type{\symbol{markcirclecross}} \NC \symbol{markcirclecross} \NC\NR
\NC \type{\symbol{markcircleheavycheck}} \NC \symbol{markcircleheavycheck} \NC \type{\symbol{markcircleheavycross}} \NC \symbol{markcircleheavycross} \NC\NR
\NC \type{\symbol{marknotcheck}} \NC \symbol{marknotcheck} \NC \NC \NC\NR
\HL
\stoptabulate

\section{Circled Text}

The \type{\circledtext} command creates numbers or text inside circles or squares.
Various styles are available through conversion commands.

\starttabulate[|l|l|p|]
\HL
\NC Parameter \NC Values \NC Description \NC\NR
\HL
\NC radius \NC 0--0.5 \NC 0=square, 0.5=circle \NC\NR
\NC framecolor \NC color \NC Border color \NC\NR
\NC fillcolor \NC color \NC Fill color \NC\NR
\NC color \NC color \NC Text color \NC\NR
\NC style \NC font style \NC Text style \NC\NR
\NC scaled \NC number \NC Scale factor \NC\NR
\HL
\stoptabulate

\startbuffer
\circledtext{1}\quad
\circledtext[radius=0.5]{A}\quad
\circledtext[radius=0,fillcolor=black,color=white]{B}\quad
\circledtext[radius=0.2]{C}

\startitemize[Co:num,packed]
  \item First item
  \item Second item
  \item Third item
\stopitemize
\stopbuffer

\typebuffer

\getbuffer

\subsection{Circled Conversions}

Available conversion sets for circled numbers:

\starttabulate[|l|l|]
\HL
\NC Conversion \NC Description \NC\NR
\HL
\NC Co:num, CO:num \NC Circled/white-on-black numbers \NC\NR
\NC Cr:num, CR:num \NC Rounded square numbers \NC\NR
\NC Cs:num, CS:num \NC Square numbers \NC\NR
\NC Co:Alph, Co:alph \NC Circled letters \NC\NR
\NC Co:cn \NC Circled Chinese numerals \NC\NR
\NC Co:hira, Co:kata \NC Circled Japanese characters \NC\NR
\HL
\stoptabulate

\chapter{More Test}

\question[I][answer=1,showanswer=true,showpoint=true]
  {one nest\answer{explain}}

\question[I][answer=1,showanswer=true,showpoint=true]{one nest
  \question[A][answer=2]{two nest
    \question[I][answer=3]{three nest\answer{explain in question}}
    \question[I][answer=4]{three nest\answer{explain in question}}
    \question[I][answer=5]{three nest\answer{explain in question}}
  }
  \question[n][answer=2]{two nest
    \question[I][answer=6]{three nest\answer{explain in question}}
    \question[I][answer=7]{three nest\answer{explain in question}}
    \question[I][answer=8]{three nest\answer{explain in question}}
  }
}

\question[showanswer=true,showpoint=true]{new test
  \problem{
    \pitem[answer=A]{new test 1\answer{explain}}
    \pitem[answer=B]{new test 2\answer{explain}}
    \pitem[answer=C]{new test 3\answer{explain}}
  }
}

\question[point=1,showpoint=true,answer=9]
{It should be reset back to 1 here.\choice{XXX,YYY,ZZZ}}
\startanswer[showanswer=true,showpoint=true]
   some text
\stopanswer

\startquestion[point={1.5},showanswer=true,showpoint=true]
  \startproblem
    \dorecurse{5}{\pitem[answer=123] {some texts}}
  \stopproblem
\stopquestion

\setupwriting[showanswer=true,showpoint=true,answer=xxx]
\setupsubwriting[answer=xxnx]
\startwriting[point=45]
  Explanation of the topic of the composition
  \startanswer
  Reference composition 1
  \stopanswer
\stopwriting
\startwriting
  Explanation of the topic of the composition
  \startsubwriting[point=450]
    Sub Composition Topic Description
    \startanswer
    Reference composition 2
    \stopanswer
  \stopsubwriting
  \startsubwriting[point=4588]
    Sub Composition Topic Description
    \startanswer
    Reference composition 3
    \stopanswer
  \stopsubwriting
\stopwriting


\stoptext


## 统计信息完整汇总

### 一、答案树结构信息

每个题目节点包含以下信息：

| 属性 | 说明 | 获取方式 |
|------|------|----------|
| `id` | 节点唯一标识（如 `jobname-1-Q-1`） | `\UnlimitedCurrentID` |
| `order` | 层级顺序（如 `1`, `2-1`, `2-2-1`） | `\UnlimitedGetAnswer{ID}{order}` |
| `answer` | 答案内容 | `\UnlimitedGetAnswer{ID}{answer}` |
| `point` | 当前节点分数（叶子节点有效） | `\UnlimitedGetAnswer{ID}{point}` |
| `totalpoint` | 该节点及子节点总分 | `\UnlimitedGetAnswer{ID}{totalpoint}` |
| `totalnumber` | 该节点及子节点总数 | `\UnlimitedGetAnswer{ID}{totalnumber}` |
| `is_leaf` | 是否为叶子节点 | `\UnlimitedGetAnswer{ID}{is_leaf}` |
| `depth` | 节点深度（0=根节点） | `\UnlimitedGetDepth` |

### 二、Structure 层级统计

| 命令 | 说明 |
|------|------|
| `\UnlimitedStructureTotalNumber{level}{number}` | 指定层级的顶层题目数量 |
| `\UnlimitedStructureTotalPoint{level}{number}` | 指定层级的总分 |
| `\UnlimitedStructureLeafNumber{level}{number}` | 指定层级的叶子节点数 |
| `\UnlimitedCurrentStructureTotalNumber{level}` | 当前层级的顶层题目数量 |
| `\UnlimitedCurrentStructureTotalPoint{level}` | 当前层级的总分 |
| `\UnlimitedCurrentStructureLeafNumber{level}` | 当前层级的叶子节点数 |

**参数说明**：
- `level`: 1=chapter, 2=section, 3=subsection...
- `number`: 该层级的编号

### 三、节点信息获取

| 命令 | 说明 |
|------|------|
| `\UnlimitedCurrentID` | 当前节点 ID |
| `\UnlimitedGetDepth` | 当前节点深度 |
| `\UnlimitedGetAnswer{ID}{key}` | 获取指定节点的属性 |
| `\UnlimitedGetCurrentExtraData{key}` | 获取当前节点的额外数据 |

**key 可选值**：`point`, `answer`, `totalpoint`, `totalnumber`, `is_leaf`, `order`

### 四、额外数据

通过 `\UnlimitedSetExtraData{key}{value}` 设置的额外信息，如：
- `totalitemnumber`: choice 环境的选项数量
- `availablecolumns`: 可用列数
- `itemmaxwidth`: 选项最大宽度

### 五、使用示例

```tex
% 获取第一章统计
第一章题目总数: \UnlimitedStructureTotalNumber{1}{1}
第一章总分: \UnlimitedStructureTotalPoint{1}{1}

% 获取指定题目信息
题目1分数: \UnlimitedGetAnswer{jobname-1-Q-1}{point}
题目1答案: \UnlimitedGetAnswer{jobname-1-Q-1}{answer}

% 在题目内部获取当前信息
\question[point=5,answer=A]{题目
  当前ID: \UnlimitedCurrentID
  当前深度: \UnlimitedGetDepth
}
```